\documentclass[10pt,a4paper,twocolumn]{article}
\usepackage[top=0.6in,bottom=0.6in,left=0.55in,right=0.55in]{geometry}
\usepackage{amsmath,amssymb,amsfonts,bm}
\usepackage{booktabs}
\usepackage{tcolorbox}
\usepackage{tabularx}

\setlength{\columnsep}{18pt}
\setlength{\parindent}{0pt}
\setlength{\parskip}{4pt}
\pagestyle{empty}

\begin{document}

\begin{center}
    {\large \textbf{NATIONAL INSTITUTE OF TECHNOLOGY CALICUT}}\\[2pt]
    {\normalsize \textbf{Department of Physics \quad $\vert$ \quad PH7302 Condensed Matter Physics (2026--27)}}\\[2pt]
    {\bfseries \scshape Mid-Semester Evaluation Rubric \& Comprehensive Solutions Manual}
\end{center}

\vspace{-0.5em}
\hrule
\vspace{0.8em}

\section*{Evaluation Architecture for the TA}
\begin{itemize}
    \item \textbf{Section A (MCQs):} Q1--Q5 are compulsory ($5 \times 1 = 5\text{ marks}$). Award full credit (1.0) for the correct option; zero otherwise.
    \item \textbf{Section B (Analytical/Lattice):} Attempt any 3 of Q6--Q9 ($3 \times 8 = 24\text{ marks}$). If all 4 are attempted, evaluate all and award marks for the best 3.
    \item \textbf{Section C (Thermal/Transport):} Q10 is compulsory (8 marks). Evaluate any 1 of Q11 or Q12 ($2 \times 8 = 16\text{ marks}$).
    \item \textbf{Total Scaled Marks:} 45. Give partial credit exactly as demarcated in the step-by-step mark rubrics below. Do not deduct marks for alternate, mathematically sound notations.
\end{itemize}

\vspace{-0.3em}
\hrule
\vspace{0.6em}

\section*{SECTION A: MCQs (1 Mark Each)}

\textbf{Q1.} Plane intercepts: $4a, 2b, \infty c$.\\[2pt]
Reciprocals of intercepts:
$$\left(\frac{1}{4}, \, \frac{1}{2}, \, \frac{1}{\infty}\right) = \left(\frac{1}{4}, \, \frac{1}{2}, \, 0\right)$$
Multiply by the lowest common denominator ($4$):
$$\left(1, \, 2, \, 0\right)$$
\textbf{Correct Option: (c)}\hfill [1.0 Mark]

\vspace{0.4em}
\textbf{Q2.} Packing fraction of FCC crystal structure.\\[2pt]
For hard spheres of radius $R$ in a cubic unit cell with lattice constant $a = 2\sqrt{2}R$ containing $Z = 4$ atoms:
$$\text{APF} = \frac{4 \times \frac{4}{3}\pi R^3}{a^3} = \frac{\frac{16}{3}\pi R^3}{16\sqrt{2}R^3} = \frac{\pi}{3\sqrt{2}} \approx 0.74$$
\textbf{Correct Option: (c)}\hfill [1.0 Mark]

\vspace{0.4em}
\textbf{Q3.} Monoatomic chain group velocity at Brillouin zone boundary.\\[2pt]
Dispersion relation: $\omega(k) = 2\sqrt{C/M}\sin(ka/2)$.\\
Group velocity:
$$v_g = \frac{d\omega}{dk} = a\sqrt{\frac{C}{M}}\cos\left(\frac{ka}{2}\right)$$
At the zone boundary $k = \pi/a$:
$$v_g\left(\frac{\pi}{a}\right) = a\sqrt{\frac{C}{M}}\cos\left(\frac{\pi}{2}\right) = 0$$
\textbf{Correct Option: (c)}\hfill [1.0 Mark]

\vspace{0.4em}
\textbf{Q4.} Dulong-Petit high-$T$ molar heat capacity.\\[2pt]
Equipartition theorem for 1 mole ($N_A$ atoms) oscillating in 3 spatial dimensions:
$$U = 3 N_A k_B T = 3 R T \implies C_V = \left(\frac{\partial U}{\partial T}\right)_V = 3R$$
\textbf{Correct Option: (a)}\hfill [1.0 Mark]

\vspace{0.4em}
\textbf{Q5.} Drude conductivity recovered from the quantum picture.\\[2pt]
Due to Pauli exclusion, interior states are blocked. Only electrons in an active window of width $\sim k_B T$ around $E_F$ can gain net momentum, moving at the Fermi velocity $v_F$. The integration over the shifted Fermi surface reproduces $\sigma = ne^2\tau/m$.\\[2pt]
\textbf{Correct Option: (b)}\hfill [1.0 Mark]

\vspace{-0.3em}
\hrule
\vspace{0.6em}

\section*{SECTION B (Attempt Any 3)}

\subsection*{Question 6 [Total: 8 Marks]}
\textbf{(a) Reciprocal Lattice Vectors Definition [2 Marks]}
\begin{itemize}
    \item \textbf{Condition [1 Mark]:} $\bm{a}_i \cdot \bm{b}_j = 2\pi\delta_{ij}$ (or $\exp(i\bm{G}\cdot\bm{R}) = 1$).
    \item \textbf{Explicit Formulas [1 Mark]:}
    $$\bm{b}_1 = 2\pi \frac{\bm{a}_2 \times \bm{a}_3}{\bm{a}_1 \cdot (\bm{a}_2 \times \bm{a}_3)}, \quad \bm{b}_2 = 2\pi \frac{\bm{a}_3 \times \bm{a}_1}{\bm{a}_1 \cdot (\bm{a}_2 \times \bm{a}_3)},$$
    $$\bm{b}_3 = 2\pi \frac{\bm{a}_1 \times \bm{a}_2}{\bm{a}_1 \cdot (\bm{a}_2 \times \bm{a}_3)}$$
\end{itemize}

\textbf{(b) Orthogonality and Interplanar Spacing [3 Marks]}
\begin{itemize}
    \item \textbf{Proof that $\bm{G}_{hkl} \perp (hkl)$ [1.5 Marks]:}\\
    Intercepts of plane $(hkl)$ on axes $\bm{a}_1, \bm{a}_2, \bm{a}_3$ are:
    $$\bm{A} = \frac{\bm{a}_1}{h}, \quad \bm{B} = \frac{\bm{a}_2}{k}, \quad \bm{C} = \frac{\bm{a}_3}{l}$$
    Vectors lying in the plane include:
    $$\bm{u} = \bm{B} - \bm{A} = \frac{\bm{a}_2}{k} - \frac{\bm{a}_1}{h}, \quad \bm{v} = \bm{C} - \bm{B} = \frac{\bm{a}_3}{l} - \frac{\bm{a}_2}{k}$$
    Taking dot product with $\bm{G}_{hkl} = h\bm{b}_1 + k\bm{b}_2 + l\bm{b}_3$:
    $$\bm{G}_{hkl} \cdot \bm{u} = \bm{G} \cdot \left(\frac{\bm{a}_2}{k} - \frac{\bm{a}_1}{h}\right) = \frac{k(2\pi)}{k} - \frac{h(2\pi)}{h} = 0$$
    Similarly, $\bm{G}_{hkl} \cdot \bm{v} = 0$. Because $\bm{G}_{hkl}$ is orthogonal to two linearly independent vectors in $(hkl)$, it is normal to the plane.
    \item \textbf{Derivation of $d_{hkl} = 2\pi/|\bm{G}_{hkl}|$ [1.5 Marks]:}\\
    The unit normal is $\hat{\bm{n}} = \bm{G}_{hkl}/|\bm{G}_{hkl}|$. The interplanar spacing $d_{hkl}$ is the projection of any intercept vector (e.g., $\bm{A} = \bm{a}_1/h$) onto $\hat{\bm{n}}$:
    $$d_{hkl} = \bm{A} \cdot \hat{\bm{n}} = \frac{\bm{a}_1}{h} \cdot \frac{\bm{G}_{hkl}}{|\bm{G}_{hkl}|} = \frac{\bm{a}_1 \cdot (h\bm{b}_1 + k\bm{b}_2 + l\bm{b}_3)}{h|\bm{G}_{hkl}|}$$
    $$d_{hkl} = \frac{h(2\pi)}{h|\bm{G}_{hkl}|} = \frac{2\pi}{|\bm{G}_{hkl}|}$$
\end{itemize}

\textbf{(c) Calculations for Copper (FCC) [3 Marks]}
\begin{itemize}
    \item \textbf{Calculation of $d_{220}$ [1.5 Marks]:}
    $$d_{hkl} = \frac{a}{\sqrt{h^2 + k^2 + l^2}}$$
    For $a = 3.61\text{ \AA}$, $(hkl) = (220)$:
    $$d_{220} = \frac{3.61\text{ \AA}}{\sqrt{2^2 + 2^2 + 0^2}} = \frac{3.61}{\sqrt{8}} = \frac{3.61}{2.8284} \approx 1.276\text{ \AA}$$
    \item \textbf{Bragg Glancing Angle $\theta$ [1.5 Marks]:}\\
    Bragg's Law: $2 d_{220} \sin\theta = n\lambda$ with $n=1$, $\lambda = 1.54\text{ \AA}$:
    $$\sin\theta = \frac{\lambda}{2 d_{220}} = \frac{1.54}{2 \times 1.2763} = \frac{1.54}{2.5526} \approx 0.6033$$
    $$\theta = \arcsin(0.6033) \approx 37.1^\circ \quad (\text{or } 0.648\text{ rad})$$
\end{itemize}

\subsection*{Question 7 [Total: 8 Marks]}
\textbf{(a) Structure Factor Definition \& Intensity [2 Marks]}
\begin{itemize}
    \item \textbf{Definition [1 Mark]:} $F_{hkl} = \sum_{j} f_j \exp[-2\pi i (h x_j + k y_j + l z_j)]$, representing the coherent scattering amplitude from all basis atoms in a single unit cell.
    \item \textbf{Diffraction Intensity [1 Mark]:} Intensity $I_{hkl} \propto |F_{hkl}|^2$. Destructive phase interference occurs when $|F_{hkl}|^2 = 0$, giving systematically absent reflections.
\end{itemize}

\textbf{(b) BCC Structure Factor Derivation [3 Marks]}
\begin{itemize}
    \item \textbf{Derivation [2 Marks]:}\\
    BCC basis: $(0,0,0)$ and $(1/2, 1/2, 1/2)$ with identical scattering factors $f_j = f$:
    $$F_{hkl} = f e^{0} + f e^{-2\pi i(h/2 + k/2 + l/2)} = f \left[1 + e^{-\pi i(h+k+l)}\right]$$
    (equivalently $f[1+e^{\pi i(h+k+l)}]$, since $e^{\pm i\pi N}=(-1)^N$ for integer $N$ — accept either sign convention).
    \item \textbf{Extinction Rule [1 Mark]:}\\
    Using $e^{-\pi i N} = (-1)^N$:
    $$F_{hkl} = \begin{cases} 2f, & \text{if } h+k+l \text{ is even (Observed)} \\ 0, & \text{if } h+k+l \text{ is odd (Extinct)} \end{cases}$$
\end{itemize}

\textbf{(c) FCC Structure Factor \& Allowed Reflections [3 Marks]}
\begin{itemize}
    \item \textbf{Derivation [1.5 Marks]:}\\
    FCC basis atoms: $(0,0,0),\, (1/2,1/2,0),\, (1/2,0,1/2),\, (0,1/2,1/2)$:
    $$F_{hkl} = f \left[1 + e^{-\pi i(h+k)} + e^{-\pi i(h+l)} + e^{-\pi i(k+l)}\right]$$
    If indices $h, k, l$ are \textbf{unmixed} (all even or all odd), pairwise sums $h+k, h+l, k+l$ are even $\implies F_{hkl} = f[1+1+1+1] = 4f$.\\[2pt]
    If indices are \textbf{mixed} (e.g., two odd, one even), two sums are odd and one is even $\implies F_{hkl} = f[1 - 1 - 1 + 1] = 0$.
    \item \textbf{Reflection Classification [1.5 Marks]:}\\
    Classify the given planes:
    \begin{itemize}
        \item $(100)$: Mixed $\implies$ \textbf{Extinct}
        \item $(110)$: Mixed $\implies$ \textbf{Extinct}
        \item $(111)$: All odd $\implies$ \textbf{Observed}
        \item $(200)$: All even $\implies$ \textbf{Observed}
        \item $(210)$: Mixed $\implies$ \textbf{Extinct}
        \item $(220)$: All even $\implies$ \textbf{Observed}
    \end{itemize}
    Allowed reflections: \textbf{(111), (200), and (220)}.
\end{itemize}

\subsection*{Question 8 [Total: 8 Marks]}
\textbf{(a) Monoatomic Chain Dispersion Derivation [4 Marks]}
\begin{itemize}
    \item \textbf{Equation of Motion [1.5 Marks]:}\\
    Consider harmonic nearest-neighbor restoring forces acting on atom $n$ of mass $M$ at equilibrium position $x_n = na$:
    $$M \frac{d^2 u_n}{dt^2} = C \left[(u_{n+1} - u_n) - (u_n - u_{n-1})\right]$$
    $$M \ddot{u}_n = C [u_{n+1} + u_{n-1} - 2u_n]$$
    \item \textbf{Plane Wave Substitution [1.5 Marks]:}\\
    Assume traveling wave solution $u_n(t) = A e^{i(kna - \omega t)}$:
    $$-M\omega^2 A e^{i(kna-\omega t)} = C A e^{i(kna-\omega t)} \left[e^{ika} + e^{-ika} - 2\right]$$
    $$-M\omega^2 = C [2\cos(ka) - 2] = -2C [1 - \cos(ka)]$$
    Using $1 - \cos(ka) = 2\sin^2(ka/2)$:
    $$M\omega^2 = 4C \sin^2(ka/2) \implies \omega^2 = \frac{4C}{M}\sin^2(ka/2)$$
    \item \textbf{Final Dispersion Form [1 Mark]:}
    $$\omega(k) = 2\sqrt{\frac{C}{M}}\left|\sin\left(\frac{ka}{2}\right)\right|$$
\end{itemize}

\textbf{(b) Physical Quantities [1 Mark]}
\begin{itemize}
    \item \textbf{Speed of sound [0.5 Mark]:} As $k \to 0$, $\sin(ka/2) \approx ka/2$:
    $$v_s = \lim_{k \to 0} \frac{\omega(k)}{k} = a\sqrt{\frac{C}{M}}$$
    \item \textbf{Cutoff frequency [0.5 Mark]:} Occurs at zone boundary $k = \pi/a$:
    $$\omega_{\max} = 2\sqrt{\frac{C}{M}} \quad \left(\nu_{\max} = \frac{\omega_{\max}}{2\pi}\right)$$
\end{itemize}

\textbf{(c) Numerical Computations [3 Marks]}
Given: $M = 40\text{ amu} = 40 \times 1.66 \times 10^{-27}\text{ kg} = 6.64 \times 10^{-26}\text{ kg}$, $a = 3 \times 10^{-10}\text{ m}$, $C = 50\text{ N/m}$.
\begin{itemize}
    \item \textbf{Maximum Frequency $\omega_{\max}$ [1 Mark]:}
    $$\sqrt{\frac{C}{M}} = \sqrt{\frac{50}{6.64 \times 10^{-26}}} = \sqrt{7.530 \times 10^{26}} \approx 2.744 \times 10^{13}\text{ rad/s}$$
    $$\omega_{\max} = 2\sqrt{\frac{C}{M}} \approx 5.488 \times 10^{13}\text{ rad/s}$$
    Converting to ordinary (cyclic) frequency, $\nu_{\max} = \omega_{\max}/2\pi$:
    $$\nu_{\max} = \frac{5.488 \times 10^{13}}{2\pi} \approx 8.73 \times 10^{12}\text{ Hz} \;(\approx 8.73\text{ THz})$$
    \item \textbf{Speed of Sound $v_s$ [1 Mark]:}
    $$v_s = a\sqrt{\frac{C}{M}} = (3 \times 10^{-10}\text{ m})(2.744 \times 10^{13}\text{ s}^{-1}) \approx 8232\text{ m/s}$$
    \item \textbf{Group Velocity at $k = 0.3\pi/a$ [1 Mark]:}
    $$v_g(k) = \frac{d\omega}{dk} = a\sqrt{\frac{C}{M}}\cos\left(\frac{ka}{2}\right) = v_s \cos\left(\frac{0.3\pi}{2}\right)$$
    $$\cos(0.15\pi) = \cos(27^\circ) \approx 0.8910$$
    $$v_g = 8232 \times 0.8910 \approx 7335\text{ m/s}$$
\end{itemize}

\subsection*{Question 9 [Total: 8 Marks]}
\textbf{(a) Diatomic Lattice Secular Determinant [3 Marks]}
\begin{itemize}
    \item \textbf{Equations of Motion [1 Mark]:}\\
    Unit cell of length $a$ with atoms $M$ at displacement $u_n$ and $m$ at $v_n$:
    $$M \ddot{u}_n = C [v_n + v_{n-1} - 2u_n]$$
    $$m \ddot{v}_n = C [u_{n+1} + u_n - 2v_n]$$
    \item \textbf{Secular Determinant [1 Mark]:}\\
    Substitute $u_n = u e^{i(kna - \omega t)}$ and $v_n = v e^{i(kna - \omega t)}$:
    $$\begin{pmatrix} 2C - M\omega^2 & -C(1 + e^{-ika}) \\ -C(1 + e^{ika}) & 2C - m\omega^2 \end{pmatrix} \begin{pmatrix} u \\ v \end{pmatrix} = 0$$
    Setting the determinant to zero:
    $$(2C - M\omega^2)(2C - m\omega^2) - C^2|1 + e^{-ika}|^2 = 0$$
    Using $|1 + e^{-ika}|^2 = 2 + 2\cos(ka) = 4\cos^2(ka/2)$, and noting $4\cos^2(ka/2) = 4[1-\sin^2(ka/2)]$, expansion of the determinant gives (\emph{corrected}: the middle term is quadratic in $\omega^2$, not $\omega^4$):
    $$Mm\,\omega^4 - 2C(M+m)\,\omega^2 + 4C^2\sin^2(ka/2) = 0$$
    \item \textbf{Dispersion Solution [1 Mark]:}\\
    Solving this quadratic in $\omega^2$ via the quadratic formula:
    $$\omega^2(k) = C\left(\frac{1}{M} + \frac{1}{m}\right) \pm C\sqrt{\left(\frac{1}{M} + \frac{1}{m}\right)^2 - \frac{4\sin^2(ka/2)}{Mm}}$$
\end{itemize}

\textbf{(b) Zone-Boundary Band Gap Derivation [2 Marks]}
At the first Brillouin zone boundary $k = \pi/a$, $\sin^2(\pi/2) = 1$:
$$\omega^2\left(\frac{\pi}{a}\right) = C\left(\frac{1}{M} + \frac{1}{m}\right) \pm C\sqrt{\left(\frac{1}{m} - \frac{1}{M}\right)^2}$$
$$= C\left(\frac{1}{M} + \frac{1}{m}\right) \pm C\left(\frac{1}{m} - \frac{1}{M}\right)$$
This decouples cleanly:
$$\omega_+^2 = \frac{2C}{m} \implies \omega_+ = \sqrt{\frac{2C}{m}} \quad \text{(Optical Branch)}$$
$$\omega_-^2 = \frac{2C}{M} \implies \omega_- = \sqrt{\frac{2C}{M}} \quad \text{(Acoustic Branch)}$$
The forbidden frequency gap is:
$$\Delta\omega = \omega_+ - \omega_- = \sqrt{\frac{2C}{m}} - \sqrt{\frac{2C}{M}}$$

\textbf{(c) Numerical Calculations [3 Marks]}
Given: $M = 40\text{ amu} = 6.64 \times 10^{-26}\text{ kg}$, $m = 20\text{ amu} = 3.32 \times 10^{-26}\text{ kg}$, $C = 50\text{ N/m}$.
\begin{itemize}
    \item \textbf{Acoustic Boundary $\omega_-(\pi/a)$ [1 Mark]:}
    $$\omega_- = \sqrt{\frac{2 \times 50}{6.64 \times 10^{-26}}} = \sqrt{1.506 \times 10^{27}} \approx 3.881 \times 10^{13}\text{ rad/s}$$
    $$\nu_- = \frac{\omega_-}{2\pi} \approx \frac{3.881\times10^{13}}{2\pi} \approx 6.18 \times 10^{12}\text{ Hz} \;(\approx 6.18\text{ THz})$$
    \item \textbf{Optical Boundary $\omega_+(\pi/a)$ [1 Mark]:}
    $$\omega_+ = \sqrt{\frac{2 \times 50}{3.32 \times 10^{-26}}} = \sqrt{3.012 \times 10^{27}} \approx 5.488 \times 10^{13}\text{ rad/s}$$
    $$\nu_+ = \frac{\omega_+}{2\pi} \approx \frac{5.488\times10^{13}}{2\pi} \approx 8.73 \times 10^{12}\text{ Hz} \;(\approx 8.73\text{ THz})$$
    \item \textbf{Band Gap $\Delta\omega$ [1 Mark]:}
    $$\Delta\omega = \omega_+ - \omega_- = (5.488 - 3.881) \times 10^{13} \approx 1.607 \times 10^{13}\text{ rad/s}$$
    $$\Delta\nu = \frac{\Delta\omega}{2\pi} \approx \frac{1.607\times10^{13}}{2\pi} \approx 2.56 \times 10^{12}\text{ Hz} \;(\approx 2.56\text{ THz})$$
\end{itemize}

\vspace{-0.3em}
\hrule
\vspace{0.6em}

\section*{SECTION C (Q10 Compulsory; Choose Q11 or Q12)}

\subsection*{Question 10 (Compulsory) [Total: 8 Marks]}
\textbf{(a) 3-D Debye Density of States Derivation [4 Marks]}
\begin{itemize}
    \item \textbf{Mode counting in $k$-space [1.5 Marks]:}\\
    Assuming 3 acoustic modes (1 longitudinal, 2 transverse) with average sound velocity $v_s$, $\omega = v_s k$. Periodic boundary conditions in volume $V$:
    $$N(k) = 3 \times \frac{\frac{4}{3}\pi k^3}{(2\pi/L)^3} = 3 \frac{V k^3}{6\pi^2} = \frac{V k^3}{2\pi^2}$$
    Expressing in terms of $\omega = v_s k$:
    $$N(\omega) = \frac{V \omega^3}{2\pi^2 v_s^3}$$
    \item \textbf{Density of states $g(\omega)$ [1.5 Marks]:}
    $$g(\omega) = \frac{dN}{d\omega} = \frac{3V \omega^2}{2\pi^2 v_s^3}$$
    \item \textbf{Debye cutoff condition [1 Mark]:}\\
    Integrating up to $\omega_D$ accounts for all $3N$ normal modes:
    $$\int_0^{\omega_D} g(\omega)\,d\omega = \frac{3V}{2\pi^2 v_s^3} \int_0^{\omega_D} \omega^2\,d\omega = \frac{V \omega_D^3}{2\pi^2 v_s^3} = 3N$$
    $$\implies \omega_D^3 = 6\pi^2 \frac{N}{V} v_s^3 \implies \omega_D = v_s \left(6\pi^2 \frac{N}{V}\right)^{1/3}$$
    Eliminating $v_s^3$ yields:
    $$g(\omega) = \frac{9N}{\omega_D^3}\omega^2 \quad (\text{for } 0 \le \omega \le \omega_D)$$
\end{itemize}

\textbf{(b) Low-$T$ Debye $T^3$ Law Derivation [1 Mark]}
Internal thermal vibrational energy:
$$U = \int_0^{\omega_D} \frac{\hbar\omega}{e^{\hbar\omega/k_B T} - 1} g(\omega)\,d\omega = \frac{9N\hbar}{\omega_D^3}\int_0^{\omega_D} \frac{\omega^3}{e^{\hbar\omega/k_B T} - 1}\,d\omega$$
Let $x = \frac{\hbar\omega}{k_B T}$ and $\Theta_D = \frac{\hbar\omega_D}{k_B}$. At low temperatures $T \ll \Theta_D$, extend upper limit $x_D = \Theta_D/T \to \infty$:
$$U = \frac{9N k_B T}{\Theta_D^3} T^3 \int_0^\infty \frac{x^3}{e^x - 1}\,dx$$
Standard integral: $\int_0^\infty \frac{x^3}{e^x - 1}dx = \frac{\pi^4}{15}$:
$$U = \frac{9Nk_B T^4}{\Theta_D^3} \left(\frac{\pi^4}{15}\right) = \frac{3\pi^4}{5} N k_B \frac{T^4}{\Theta_D^3}$$
Differentiating with respect to $T$:
$$C_V = \left(\frac{\partial U}{\partial T}\right)_V = \frac{12\pi^4}{5} N k_B \left(\frac{T}{\Theta_D}\right)^3$$

\textbf{(c) Calculations for Sodium [3 Marks]}
Given: $\nu_D = 3.3 \times 10^{12}\text{ Hz}$ (i.e., $\omega_D = 2\pi\nu_D \approx 2.074\times10^{13}\text{ rad/s}$), $T = 10\text{ K}$, $N_A = 6.022 \times 10^{23}\text{ mol}^{-1}$.
\begin{itemize}
    \item \textbf{Debye Temperature $\Theta_D$ [1.5 Marks]:}
    $$\Theta_D = \frac{h \nu_D}{k_B} = \frac{6.626 \times 10^{-34}\text{ J}\cdot\text{s} \times 3.3 \times 10^{12}\text{ s}^{-1}}{1.38 \times 10^{-23}\text{ J/K}}$$
    $$\Theta_D = \frac{2.1866 \times 10^{-21}}{1.38 \times 10^{-23}} \approx 158.45\text{ K}$$
    \item \textbf{Molar Heat Capacity at $T = 10\text{ K}$ [1.5 Marks]:}\\
    Here $T/\Theta_D = 10 / 158.45 \approx 0.0631 < 0.1$, so the $T^3$ approximation is valid.
    $$C_V = \frac{12\pi^4}{5} R \left(\frac{T}{\Theta_D}\right)^3$$
    With $R = N_A k_B = 8.314\text{ J mol}^{-1}\text{K}^{-1}$:
    $$\frac{12\pi^4}{5} R = 2.4 \times 97.409 \times 8.314 \approx 1943.6\text{ J mol}^{-1}\text{K}^{-1}$$
    $$C_V = 1943.6 \times (0.06311)^3 = 1943.6 \times (2.514 \times 10^{-4})$$
    $$C_V \approx 0.489\text{ J mol}^{-1}\text{K}^{-1}$$
\end{itemize}

\subsection*{Question 11 [Total: 8 Marks]}
\textbf{(a) Free Electron DOS \& $E_F$ Derivation [4 Marks]}
\begin{itemize}
    \item \textbf{Density of states derivation [2.5 Marks]:}\\
    In a cubic box of volume $V = L^3$, periodic boundary conditions yield allowed wavevectors $k_i = \frac{2\pi n_i}{L}$. Each mode occupies $k$-space volume $(2\pi/L)^3 = 8\pi^3/V$.\\
    Number of electron states (including spin degeneracy $g_s = 2$) with wavevector inside sphere of radius $k$:
    $$N(k) = 2 \times \frac{\frac{4}{3}\pi k^3}{(2\pi)^3/V} = \frac{V k^3}{3\pi^2}$$
    Using $E = \frac{\hbar^2 k^2}{2m} \implies k = \left(\frac{2mE}{\hbar^2}\right)^{1/2}$:
    $$N(E) = \frac{V}{3\pi^2}\left(\frac{2mE}{\hbar^2}\right)^{3/2} = \frac{V}{3\pi^2}\left(\frac{2m}{\hbar^2}\right)^{3/2} E^{3/2}$$
    Taking the derivative:
    $$D(E) = \frac{dN}{dE} = \frac{V}{2\pi^2}\left(\frac{2m}{\hbar^2}\right)^{3/2} \sqrt{E}$$
    \item \textbf{Fermi energy derivation [1.5 Marks]:}\\
    At $T = 0\text{ K}$, all $N$ electrons fill states up to $k_F$:
    $$N = \frac{V k_F^3}{3\pi^2} \implies n = \frac{N}{V} = \frac{k_F^3}{3\pi^2} \implies k_F = (3\pi^2 n)^{1/3}$$
    $$E_F = \frac{\hbar^2 k_F^2}{2m} = \frac{\hbar^2}{2m}(3\pi^2 n)^{2/3}$$
\end{itemize}

\textbf{(b) Calculations for Sodium [2 Marks]}
Given: $\rho = 970\text{ kg/m}^3$, $A = 23\text{ g/mol} = 0.023\text{ kg/mol}$, monovalent ($z=1$).
\begin{itemize}
    \item \textbf{Conduction Electron Density $n$ [0.5 Mark]:}
    $$n = \frac{\rho N_A}{A} = \frac{970 \times 6.022 \times 10^{23}}{0.023} \approx 2.54 \times 10^{28}\text{ m}^{-3}$$
    \item \textbf{Fermi Energy $E_F$ [1 Mark]:}
    $$3\pi^2 n = 3\pi^2 (2.54 \times 10^{28}) = 7.521 \times 10^{29}\text{ m}^{-3}$$
    $$(3\pi^2 n)^{1/3} \approx 9.094 \times 10^9\text{ m}^{-1}$$
    $$E_F = \frac{(1.054 \times 10^{-34})^2}{2 \times 9.11 \times 10^{-31}} (9.094 \times 10^9)^2$$
    $$E_F = \frac{1.1109 \times 10^{-68}}{1.822 \times 10^{-30}} \times 8.270 \times 10^{19} \approx 5.04 \times 10^{-19}\text{ J}$$
    In electron-volts:
    $$E_F = \frac{5.04 \times 10^{-19}\text{ J}}{1.602 \times 10^{-19}\text{ J/eV}} \approx 3.15\text{ eV}$$
    \item \textbf{Fermi Temperature $T_F$ [0.5 Mark]:}
    $$T_F = \frac{E_F}{k_B} = \frac{5.04 \times 10^{-19}\text{ J}}{1.38 \times 10^{-23}\text{ J/K}} \approx 36,520\text{ K}$$
\end{itemize}

\textbf{(c) Ground-State Energy \& Active Fraction [2 Marks]}
\begin{itemize}
    \item \textbf{Average Kinetic Energy at $T = 0\text{ K}$ [1 Mark]:}
    $$\langle E \rangle_0 = \frac{U_0}{N} = \frac{1}{N}\int_0^{E_F} E \, D(E)\, dE$$
    Since $D(E) = C\sqrt{E}$ and $N = \int_0^{E_F} C\sqrt{E}\,dE = \frac{2}{3} C E_F^{3/2}$:
    $$U_0 = C \int_0^{E_F} E^{3/2}\,dE = \frac{2}{5} C E_F^{5/2} \implies \langle E \rangle_0 = \frac{3}{5} E_F$$
    $$\langle E \rangle_0 = 0.6 \times 3.15\text{ eV} \approx 1.89\text{ eV} \quad (3.02 \times 10^{-19}\text{ J})$$
    \item \textbf{Fraction Excited at $T = 300\text{ K}$ [1 Mark]:}\\
    Thermal energy range: $\Delta E = 2 k_B T$. Conduction electrons excited around $E_F$:
    $$f_{\text{th}} \approx \frac{D(E_F) \times 2 k_B T}{N}$$
    Since $D(E_F) = \frac{3N}{2E_F}$:
    $$f_{\text{th}} \approx \frac{3}{2}\frac{2 k_B T}{E_F} = 3 \frac{k_B T}{E_F} = 3 \frac{T}{T_F}$$
    $$f_{\text{th}} \approx 3 \times \frac{300\text{ K}}{36,520\text{ K}} \approx 3 \times 0.00821 \approx 0.0246 \quad (2.5\%)$$
    \textit{(Full credit if student uses simplified scaling $f \approx T/T_F \approx 0.8\%$).}
\end{itemize}

\subsection*{Question 12 [Total: 8 Marks]}
\textbf{(a) Boltzmann Transport Derivation of $\sigma$ [4 Marks]}
\begin{itemize}
    \item \textbf{Steady-State Boltzmann Transport Equation [1 Mark]:}\\
    For spatially uniform electron gas under DC field $\bm{\mathcal{E}} = \mathcal{E}_x \hat{\bm{x}}$, the relaxation-time approximation gives:
    $$\dot{\bm{k}} \cdot \nabla_{\bm{k}} f = -\frac{f - f_0}{\tau}$$
    \item \textbf{Perturbed Distribution $f(\bm{k})$ [1.5 Marks]:}\\
    Semiclassical force: $\hbar \dot{k}_x = -e\mathcal{E}_x$. Expanding to first order ($f \approx f_0$ on LHS):
    $$\dot{k}_x \frac{\partial f_0}{\partial k_x} = \left(-\frac{e\mathcal{E}_x}{\hbar}\right)\left(\hbar v_x \frac{\partial f_0}{\partial E}\right) = -e\mathcal{E}_x v_x \frac{\partial f_0}{\partial E}$$
    $$-e\mathcal{E}_x v_x \frac{\partial f_0}{\partial E} = -\frac{f - f_0}{\tau} \implies f(\bm{k}) = f_0 + e\tau \mathcal{E}_x v_x \left(-\frac{\partial f_0}{\partial E}\right)$$
    \item \textbf{Current Density and Conductivity [1.5 Marks]:}\\
    Current density along $x$:
    $$j_x = -e \int \frac{2\, d^3 k}{8\pi^3} v_x f(\bm{k})$$
    Because $v_x f_0(E)$ is odd in $k_x$, the equilibrium integral vanishes:
    $$j_x = \frac{e^2 \tau \mathcal{E}_x}{4\pi^3} \int v_x^2 \left(-\frac{\partial f_0}{\partial E}\right) d^3 k$$
    At low temperatures, $-\frac{\partial f_0}{\partial E} \approx \delta(E - E_F)$. For cubic crystals, $\langle v_x^2 \rangle = \frac{1}{3}v_F^2$:
    $$j_x = \frac{e^2 \tau \mathcal{E}_x}{4\pi^3}\left(\frac{1}{3}v_F^2\right) \left(\frac{4\pi k_F^2}{\hbar v_F}\right) = \frac{e^2 \tau k_F^3}{3\pi^2 m}\mathcal{E}_x$$
    Using $n = k_F^3/(3\pi^2)$, this simplifies to $j_x = \frac{n e^2 \tau}{m} \mathcal{E}_x$.\\
    Comparing with Ohm's Law $j_x = \sigma \mathcal{E}_x$:
    $$\sigma = \frac{n e^2 \tau}{m}$$
\end{itemize}

\textbf{(b) Calculations for Copper [2 Marks]}
Given: $n = 8.5 \times 10^{28}\text{ m}^{-3}$, $\tau = 2.5 \times 10^{-14}\text{ s}$, $\mathcal{E} = 1\text{ V/m}$.
\begin{itemize}
    \item \textbf{Conductivity $\sigma$ [1 Mark]:}
    $$\sigma = \frac{(8.5 \times 10^{28})(1.602 \times 10^{-19})^2 (2.5 \times 10^{-14})}{9.11 \times 10^{-31}}$$
    $$\sigma = \frac{8.5 \times 10^{28} \times 2.566 \times 10^{-38} \times 2.5 \times 10^{-14}}{9.11 \times 10^{-31}}$$
    $$\sigma = \frac{5.454 \times 10^{-23}}{9.11 \times 10^{-31}} \approx 5.99 \times 10^7\text{ }\Omega^{-1}\text{m}^{-1}$$
    \item \textbf{Drift Velocity $v_d$ [1 Mark]:}
    $$v_d = \mu \mathcal{E} = \frac{e\tau}{m}\mathcal{E} = \frac{(1.602 \times 10^{-19})(2.5 \times 10^{-14})}{9.11 \times 10^{-31}} \times 1$$
    $$v_d = \frac{4.005 \times 10^{-33}}{9.11 \times 10^{-31}} \approx 4.40 \times 10^{-3}\text{ m/s} = 4.40\text{ mm/s}$$
\end{itemize}

\textbf{(c) Fermi Velocity \& Mean Free Path [2 Marks]}
\begin{itemize}
    \item \textbf{Fermi Velocity $v_F$ [0.75 Mark]:}
    $$k_F = (3\pi^2 n)^{1/3} = (3\pi^2 \times 8.5 \times 10^{28})^{1/3} \approx (2.517 \times 10^{30})^{1/3}$$
    $$k_F \approx 1.360 \times 10^{10}\text{ m}^{-1}$$
    $$v_F = \frac{\hbar k_F}{m} = \frac{(1.054 \times 10^{-34})(1.360 \times 10^{10})}{9.11 \times 10^{-31}} \approx 1.57 \times 10^6\text{ m/s}$$
    \item \textbf{Mean Free Path $l$ [0.75 Mark]:}
    $$l = v_F \tau = (1.57 \times 10^6\text{ m/s}) \times (2.5 \times 10^{-14}\text{ s})$$
    $$l \approx 3.93 \times 10^{-8}\text{ m} = 39.3\text{ nm} = 393\text{ \AA}$$
    \item \textbf{Lattice Comparison [0.5 Mark]:}\\
    Ratio relative to interatomic spacing $a = 3.61\text{ \AA}$:
    $$\frac{l}{a} = \frac{393\text{ \AA}}{3.61\text{ \AA}} \approx 109$$
    \textbf{Physical conclusion:} The electronic mean free path is more than $100$ lattice constants long. This demonstrates that electrons traverse hundreds of unit cells without scattering, because an electron wave propagates without resistance through a perfectly periodic crystal potential. Collisions occur only due to lattice defects, impurities, or thermal vibrations (phonons).
\end{itemize}

\end{document}